% =============================================================================
% Chapter 6  Coupled Systems and Complex-Valued Fields  --  ~2.75 pp
% rubric: Project Body (Modelling) -- Model Formulation /30.
% Plan: ThesisPlanning.md §5.6 (T-31, T-32; new in v2.1).
%
% PURPOSE     Take the construction from one real scalar field, the paper's scope,
%             to several coupled fields, and with them to complex-valued linear
%             equations; then hand a coupled linear system to Chapter 7, since a
%             Carleman lift is one. It answers no numbered question (T-31).
% CRITERION   Formulation, "sufficient detail to permit readers to create
%             equivalent models": the rows, the data and the circuits of a system
%             must be rebuildable from the body plus Appendix H. Opens with the
%             scope the paper leaves (one real field; its Discussion: with a
%             Fourier basis "the resulting total energy function is of a
%             complicated form, rendering it difficult to find the ground state")
%             and closes on Q(iii)'s prerequisite.
% SOURCES     ThesisPlanning.md §5.6, §3.4, §8; pihm/docs/02_methods.md §3.2 (the
%             mass form), §3.3 (ratio rows), §3.4, §6.2 (a system's initial state);
%             pihm/docs/companion/{21,51,61}; pihm/docs/16_extensions.md §1.1 (R7,
%             the mass form against the chain), §2-4; pihm/docs/03_problems/07_r7.md;
%             decisions 39, 42, 45 (reasons only); src/pihm/assemble/systems.py,
%             circuits/systems.py, ir.py (matching_terms), operators/constraints.py
%             (ratio_rows), problems/fluids.py (real_embedding), circuits/tensor.py
%             (complex weights); notebook 60; the paper's Discussion (p. 11) for
%             the complex-valued remark.
% WRITE WHEN  After W0c step 3 has moved the systems material here (the coupled-fields
%             clause of sec:descriptor-constraints, the matchings paragraph of
%             sec:descriptor-block-encoding, the space-time discretisations and
%             ratio-row item of sec:carleman); the mathematics with numbers as \res,
%             in full with Chapter 5 at gate G4 (wave W3).
% =============================================================================

% LAND (opening, no heading; 0.2 p)
%   The paper's scope (one real scalar field), its remark on complex-valued states, quoted
%   ("the resulting total energy function is of a complicated form, rendering it difficult to
%   find the ground state"), and what this chapter adds; the roadmap.
% REVISE (v2.1): the opening carries no claim beyond scope; it answers no numbered question (T-31).

\subsection{Several Fields on One State}
\label{sec:systems-fields}
% ~0.6 pp
% LAND, in order:
%   1. A first-order system u_i' = sum_j A_ij u_j (constant coefficients; one equation per field,
%      each holding its own first derivative) on one Chebyshev axis; on a time axis xi = 2t/T - 1,
%      with the factor T/2 (moved here from Ch. 7's space-time paragraph).
%   2. The state: one block per component, on a field register of ceil(log2 C) qubits, padded to a
%      power of two with padding rows of weight 1 that decouple exactly.
%   3. The rows: standard I (x) G - A (x) I (the paper's elimination, applied per field); the
%      descriptor's mass form (sec:descriptor-layouts) multiplies them by I (x) B~. Same kernel; an
%      IVP with data has a one-dimensional near-kernel.
%   4. Higher-order systems: by reduction to first order (one sentence; none run beyond R7's
%      order-two equation written as two fields).
%   MOVED VERBATIM (W0c step 3):
%     (a) from 5_descriptor.tex, sec:descriptor-constraints (the coupled-fields clause):
%       A system of coupled first-order fields takes the mass form, one block per component.
%     (c) from 7_nonlinear.tex, sec:carleman, LAND item 3 (the two space-time discretisations):
%       3. The two time discretisations (eq:carleman-spacetime): the same kernel,
%          differing only by a banded invertible left factor -- why R8-R10 compare
%          like for like. The descriptor's time axis lowers the degree.
%       (eq:carleman-spacetime itself stays in Ch. 7; its factor T/2 and axis map are stated here.)
% OFFLOAD: the system rows in full and the padding -> Appendix H (app:systems).
% TRAPS: redefining the mass form (sec:descriptor-layouts owns it); claiming higher-order systems
%   were run; "vector-valued solver".

\subsection{Data for Several Fields: Ratio Rows}
\label{sec:systems-data}
% ~0.4 pp
% LAND, in order:
%   1. The ratio row F = y_a (I - |a><a|) - w <a| on the field register, times B(x*): it imposes the
%      direction of the vector datum and leaves its scale free, as a zero constraint does; divided
%      by y_a it is scale-free; a zero datum holds its component to zero exactly.
%   2. The scale from the regular datum, as for one field. The initial state's weighting by
%      w0/||w0|| is stated in sec:pipeline (plan §8.5); one pointer, with the reason (equal weights
%      cancel on Taylor-Green).
%   MOVED VERBATIM (W0c step 3), from 7_nonlinear.tex, sec:carleman, LAND item 4 (ratio-row data):
%       4. Initial data enter through ratio rows, F (x) <tau(-1)| (tab:constraint-kinds).
% OFFLOAD: the ratio rows and their three-term circuit -> Appendix H (app:systems).
% TRAPS: a ratio row as a fixed-scale constraint; restating the weighting rule.

\subsection{The Mass Form for Systems}
\label{sec:systems-mass}
% ~0.4 pp
% LAND, in order:
%   1. Why the mass form and not a chain per field: the equation already gives each derivative, so a
%      derivative block carries nothing new, costs register and needs rescaling. Measured on R7
%      (u' = v, v' = -pi^2 u), the mass form's gap is about 2.1 times the block-rescaled chain's,
%      flat in n (key to add, plan §2.1); the chain stays as a labelled ablation.
%   2. The standard and descriptor forms of a system share their state, initial state and readout,
%      and differ only by I (x) B~: the like-for-like comparison R7-R10 rely on.
% OFFLOAD: none.
% TRAPS: redefining the mass form or restating its kernel theorem (sec:descriptor-layouts owns both:
%   one theorem covers both layouts); "about 2.1x" without the R7 and block-rescaled-chain scope.

\subsection{The Coupling on a Circuit}
\label{sec:systems-coupling}
% ~0.4 pp
% LAND, in order:
%   1. A split by edge colouring into as many matchings as its largest row or column degree; each a
%      permutation of the field register's basis states, its entries a multiplexed rotation on one
%      flag; alpha = the largest entry. Both forms pay the same for A.
%   2. The ratio rows as three LCU terms (identity, projector, a loaded rank-one term).
%   3. The matchings against Pauli words, one sentence with numbers (key to add); construction ->
%      Appendix H.
%   MOVED VERBATIM (W0c step 3):
%     (b) from 5_descriptor.tex, sec:descriptor-block-encoding (the matchings note):
%       % MOVED 2026-10-06 to 6_nonlinear.tex (the Carleman coupling's matchings, where the lift is built):
%       %   a large irregular coupling is split into matchings, each with at most one entry in every row and
%       %   column, and a greedy colouring uses at most 2*Delta - 1 of them, Delta the largest row degree.
%     (d) from 7_nonlinear.tex, sec:folds, LAND item 4 (the symmetric layout's matchings):
%       "as matchings in the symmetric layout"
%       MOVED here from 5_descriptor.tex (2026-10-06): a large irregular coupling is not expanded in Pauli
%       words but split into matchings, each with at most one entry in every row and column; a greedy
%       colouring uses at most 2*Delta - 1 of them, Delta the largest row degree.
%     (The tensor layout's 2K - 1 terms stay in Ch. 7, sec:folds.)
% REVISE (v2.1): (b) and (d) are the same sentence moved twice; keep one in the drafted prose. The
%   colouring bound (2*Delta - 1, greedy) and the plan's "as many matchings as the largest row or
%   column degree" (edge colouring, Delta) differ; check against circuits/systems.py and ir.py
%   (matching_terms) before drafting.
% OFFLOAD: the matchings by edge colouring -> Appendix H (app:systems).
% TRAPS: calling the matchings a Pauli expansion; claiming a gate count the key does not give.

\subsection{Complex-Valued Fields}
\label{sec:systems-complex}
% ~0.75 pp. Scope fixed by T-32 (do not widen).
% LAND, in order:
%   1. The formulation does not need a non-Hermitian filter. For complex coefficients, data or
%      sources R is complex, and H = R^dagger R is still Hermitian and positive semi-definite; the
%      ground state is still the smallest right singular vector of R; the stacked reflection and
%      GQSP act on H' unchanged. The non-Hermitian operator is R, which the filter never sees.
%      (Narrative and mathematical; no complex-R preparation was run.)
%   2. Realification, the route run (T-32): z = u + iv as two real fields; z' = cz gives
%      [u; v]' = [[Re c, -Im c], [Im c, Re c]] [u; v] (c, a complex coefficient, never a or alpha,
%      which are taken); one more field-register qubit; the same kernel under the real-linear
%      isomorphism. R7 is exactly this: with z = pi u + iv, u' = v and v' = -pi^2 u are
%      z' = -i pi z, a single-mode Schroedinger-type equation (R7 = z' = -i pi z realified). The
%      Fourier-Galerkin fluids of Ch. 7 are carried on real bases for the same reason (sine modes
%      for Burgers; cos/sin with an isometric real embedding for Navier-Stokes). R9 and R10 use
%      real Fourier bases.
%   3. Native complex encoding, built once (T-32): the tensor layout's coupling (Ch. 7) is built on
%      complex Fourier variables, with complex phases i^{sum a} in its LCU weights, and verified
%      (exact and IR); it is never used to prepare a state, because the classical pipeline carries
%      real data only. It is the one native complex encoding: verified, never prepared.
%   4. Not built: complex sources, data and coefficients in the circuits' loaders (a phase on the
%      collapse row; a complex state loader), and so native complex preparation. `pihm` refuses
%      complex sources, data and coefficients. Future work (§11.2).
%   5. The research proposal's "complex-valued, non-Hermitian equations" is answered here at the
%      level of the formulation, and recorded as a deviation (Appendix M, plan §13.4).
% OFFLOAD: the realification map and its kernel equivalence, the real embedding of the Fourier
%   variables -> Appendix H (app:systems).
% TRAPS: "handles complex equations" or "complex solver" unscoped; "a complex-valued field was
%   prepared or solved"; calling the tensor encoding "prepared"; the legacy DESolverLib's complex
%   code is NEVER cited (any complex-valued result resting on it is out); a non-Hermitian filter
%   (none is needed, and none is run).

\subsection{From Systems to Lifts}
\label{sec:systems-lifts}
% ~0.2 pp. Hand-off to Ch. 7.
% LAND, in order:
%   1. A Carleman lift of a polynomial system is a linear system on its monomials, so every
%      construction of this chapter carries over: the components become monomials, the ratio rows
%      carry the data's moments, and A becomes A_K. Ch. 7 adds only the lift and its products.
% OFFLOAD: none.
% TRAPS: re-deriving the lift (Ch. 7 owns it); stating the lift prepares without the convergence
%   caveat (rho < 1 is sufficient, convergence is measured).

% CHECKLIST (marker): the scope gap stated from the paper's own words; systems rebuildable from body
%   and Appendix H; the mass form applied, not redefined (sec:descriptor-layouts owns it);
%   realification stated as the route run, with R7 as its instance; every complex-valued claim
%   scoped (formulation; one encoding verified; preparation not built); the hand-off to Ch. 7.
% TRAPS (chapter): "handles complex equations" unscoped; any complex-valued result resting on the
%   legacy DESolverLib (never cited); redefining the mass form or restating its kernel theorem;
%   calling the tensor encoding "prepared"; claiming higher-order systems were run;
%   "vector-valued solver".
